%!TEX root = ../main.tex

\section{Introduction}\label{sec:introduction}

Strategy Incorporated (ticker: MSTR) holds 761,068 Bitcoin---approximately
3.8\% of the circulating supply---valued at \$56.25 billion as of March~2026.
Its equity trades at a persistent premium to net asset value: the market
capitalisation exceeds the mark-to-market value of the BTC holdings minus
liabilities.  This premium, which has ranged from near zero to over 200\%,
represents one of the most striking pricing anomalies in modern corporate
finance.  Yet no structural theory explains \emph{what} the fair premium
should be, \emph{when} it is justified, or \emph{how} it interacts with the
firm's capital structure decisions.

This paper fills the gap.  We develop a theory in which the NAV premium
is endogenous: it equals the present value of a \emph{perpetual accumulation
option}---the right to issue equity at a premium to NAV and deploy the
proceeds to purchase BTC, thereby increasing BTC-per-share for existing
holders.  The key insight is that this option is self-referential: the premium
that justifies the option is itself sustained by the option's value.  This
reflexive structure generates rich equilibrium dynamics, including multiple
stable equilibria and path-dependent transitions.

\paragraph{Contributions.}
The paper makes five contributions to the literature on alternative-asset
holding companies and crypto-finance:

\begin{enumerate}[leftmargin=*]
  \item \textbf{Endogenous premium theory with reflexivity.}
    We derive a fixed-point equation for the fair premium
    $\pi^\star$ that links it to BTC volatility, issuance intensity,
    leverage, and discount rates (Proposition~\ref{prop:fair_premium},
    Proposition~\ref{prop:steady_state}).  The closed-form steady-state
    approximation connects to the Black--Scholes call-option formula,
    providing both analytical tractability and economic intuition.

  \item \textbf{Multiple equilibria and death spiral.}
    The reflexive feedback loop creates three equilibria: a high-premium
    ``flywheel'' (virtuous cycle of issuance and accumulation), an unstable
    tipping point, and a low-premium spiral
    (Theorem~\ref{thm:three_equilibria}).  Transitions between equilibria
    exhibit hysteresis via saddle-node bifurcations
    (Proposition~\ref{prop:hysteresis}), and we characterise the death
    spiral mechanism formally (Theorem~\ref{thm:death_spiral}).

  \item \textbf{Greeks reinterpretation.}
    We show that five structural indicators---Implied Leverage Elasticity
    (ILE), Total Equity Elasticity (TEE), Premium Mean-Reversion Index
    (PMRI), Implied BTC Growth Rate (IBGR), and Implied Funding
    Requirement Distribution (IFRD)---map precisely onto the option Greeks
    of the embedded accumulation claim: Delta, Adjusted Delta, Theta,
    Moneyness, and Exercise Cost (Table~\ref{tab:greek-map}).  We further
    introduce Vega and Rho as new sensitivity measures.

  \item \textbf{Optimal capital structure.}
    We derive a premium-dependent financing rule that determines the
    cost-minimising channel among ATM equity, convertible debt, and
    preferred stock as a function of the log-premium
    (Proposition~\ref{prop:optimal_channel}).  We introduce the Weighted
    Average Cost of BTC Acquisition (WACBA) and characterise endogenous
    procyclical leverage dynamics (Proposition~\ref{prop:procyclical}).
    Strategy's observed financing behaviour---shifting from equity to
    convertibles to preferred stock as the premium declined through
    2025--2026---closely matches the model's predictions.

  \item \textbf{Empirical calibration through 2026.}
    We calibrate the model to a comprehensive daily panel of BTC prices,
    MSTR equity data, holdings, debt, and five series of preferred stock
    (STRK, STRF, STRC, STRD, STRE).  The calibrated premium process has
    mean-reversion speed $\hat\kappa = 5.21$ and long-run mean
    $\hat\theta = 1.08$, with BTC volatility of 48.6\%.  The current
    premium of 24.2\% sits well below the long-run mean
    ($\text{PMRI} = -0.71$), implying positive carry.
    Monte Carlo simulation with 5,000 paths yields a 91.9\% three-year
    survival probability, mean implied funding requirement of \$35.3B, and
    preferred dividend coverage ratio of $49\times$ currently.
\end{enumerate}

\paragraph{Broader implications.}
The framework applies beyond MSTR to any corporate treasury vehicle that
holds a volatile asset and issues securities to fund accumulation---including
the growing cohort of public companies adding BTC to their balance sheets.
Our market impact analysis (Section~\ref{sec:market_impact_theory})
formalises the systemic risk from corporate BTC concentration and derives
conditions for contagion cascades among leveraged holders.

\paragraph{Outline.}
Section~\ref{sec:literature} reviews related literature.
Section~\ref{sec:theory} develops the theoretical framework:
the accumulation option (Section~\ref{sec:fair_premium}),
reflexivity and multiple equilibria (Section~\ref{sec:multiple_equilibria}),
indicators as Greeks (Section~\ref{sec:greeks}),
MSTR vs.\ BTC ETF comparison (Section~\ref{sec:etf_comparison}),
optimal capital structure (Section~\ref{sec:optimal_capital}),
and market impact (Section~\ref{sec:market_impact_theory}).
Section~\ref{sec:empirical} describes the data and estimation.
Section~\ref{sec:results} presents calibrated results.
Section~\ref{sec:discussion} discusses implications and limitations.
Section~\ref{sec:conclusion} concludes.
